\section*{Chapitre \ref{ch:g2:mixing-and-dissolving} --- Mélanger et dissoudre}

\begin{solution}{exo:g2:mixing-and-dissolving:1}
Le sucre et le sel se dissolvent. Le sable, la craie et les cailloux ne
se dissolvent pas.
\end{solution}

\begin{solution}{exo:g2:mixing-and-dissolving:2}
Une \omterm{def:g2:mixing-and-dissolving:dissolve}{solution} est le liquide limpide obtenu quand un solide \omterm{def:g2:mixing-and-dissolving:dissolve}{se dissout}
dans l'eau. Exemples~: l'eau sucrée, l'eau salée (le thé sucré en est
une aussi).
\end{solution}

\begin{solution}{exo:g2:mixing-and-dissolving:3}
Oui~: on a mis ensemble des flocons d'avoine, des raisins secs, des noix
et des graines. On voit encore chaque partie et on peut la reprendre.
\end{solution}

\begin{solution}{exo:g2:mixing-and-dissolving:4}
Non. Un solide dissous laisse l'eau limpide, sans rien au fond~; ici,
l'eau est trouble et une couche s'est déposée.
\end{solution}

\begin{solution}{exo:g2:mixing-and-dissolving:5}
Il est limpide (on voit à travers) mais pas incolore (il est rose).
C'est une \omterm{def:g2:mixing-and-dissolving:dissolve}{solution}~: une \omterm{def:g2:mixing-and-dissolving:dissolve}{solution} peut avoir une couleur.
\end{solution}

\begin{solution}{exo:g2:mixing-and-dissolving:6}
Le thé a un goût sucré~: le sucre est toujours dedans, réparti dans
l'eau en morceaux trop petits pour être vus.
\end{solution}

\begin{solution}{exo:g2:mixing-and-dissolving:7}
Le soleil fait sécher l'eau~; le sel qui était dissous dans l'eau de mer
reste.
\end{solution}

\begin{solution}{exo:g2:mixing-and-dissolving:8}
«~Dissous.~» Fondre, c'est ce que fait un glaçon tout seul, sans eau~;
le sucre a besoin de l'eau, et il est toujours dans le thé.
\end{solution}

\begin{solution}{exo:g2:mixing-and-dissolving:9}
Le bécher qui contient le sable est trouble et a une couche au fond~;
l'eau sucrée est limpide, sans rien au fond.
\end{solution}

\begin{solution}{exo:g2:mixing-and-dissolving:10}
L'eau du robinet n'est pas seulement de l'eau~: quelque chose y est
dissous. Quand l'eau sèche, ce solide reste~: c'est le cercle blanc.
\end{solution}

\begin{solution}{exo:g2:mixing-and-dissolving:11}
Elle laisse sécher quelques gouttes de chacun dans une coupelle sombre
et propre. L'eau pure ne laisse rien~; l'eau salée laisse du sel blanc.
\end{solution}
